
\documentclass{article}
\usepackage{colortbl}
\usepackage{makecell}
\usepackage{multirow}
\usepackage{supertabular}

\begin{document}

\newcounter{utterance}

\twocolumn

{ \footnotesize  \setcounter{utterance}{1}
\setlength{\tabcolsep}{0pt}
\begin{supertabular}{c@{$\;$}|p{.15\linewidth}@{}p{.15\linewidth}p{.15\linewidth}p{.15\linewidth}p{.15\linewidth}p{.15\linewidth}}

    \# & $\;$A & \multicolumn{4}{c}{Game Master} & $\;\:$B\\
    \hline 

    \theutterance \stepcounter{utterance}  

    & & \multicolumn{4}{p{0.6\linewidth}}{\cellcolor[rgb]{0.9,0.9,0.9}{%
	\makecell[{{p{\linewidth}}}]{% 
	  \tt {\tiny [A$\langle$GM]}  
	 State giocando a un gioco di negoziazione cooperativa in cui dovete accordarvi su come dividere una serie di oggetti tra voi e un altro giocatore.\\ \tt \\ \tt Regole:\\ \tt (a) A Lei e all'altro giocatore viene dato un insieme di oggetti. Ognuno di voi riceve anche una funzione di valore segreta, che rappresenta il valore di ciascun tipo di oggetto.\\ \tt (b) Si scambiano messaggi con l'altro giocatore per concordare chi si aggiudica gli oggetti. Potete inviare un massimo di 5 messaggi ciascuno, oppure terminare in anticipo facendo una proposta segreta in qualsiasi momento.\\ \tt (c) A ciascuno di voi viene chiesto di inviare una proposta segreta indicando gli oggetti che desiderate, formattata tra parentesi quadre come segue: "[Proposta: <numero> <nome oggetto>, <numero> <nome oggetto>, <...>]".\\ \tt (d) Se le vostre proposte sono complementari, cioè ci sono abbastanza oggetti per soddisfare entrambe le proposte, a ciascun giocatore viene assegnato un punteggio basato sulla somma dei valori degli oggetti ricevuti. In caso contrario, entrambi ricevono zero punti.\\ \tt (e) L'obiettivo è massimizzare la somma del proprio punteggio e di quello dell'altro giocatore.\\ \tt \\ \tt Cominciamo.\\ \tt \\ \tt L'insieme degli oggetti disponibili è:\\ \tt \\ \tt 1 bullone, 1 quadro, 1 uovo, 4 schede\\ \tt \\ \tt La funzione valore segreta è:\\ \tt \\ \tt bullone: 6, quadro: 4, uovo: 0, scheda: 0\\ \tt \\ \tt IMPORTANTE: i vostri messaggi, a meno che non si tratti di una proposta segreta, vengono trasmessi direttamente all'altro giocatore, quindi non includete alcuna risposta alle regole o testo di annuncio del vostro messaggio. Per fare una proposta segreta, utilizzate il formato indicato. Non utilizzare delle parentesi quadre quando si comunica all'altro giocatore, altrimenti verrà interpretata come una proposta segreta.\\ \tt  
	  } 
	   } 
	   } 
	 & \\ 
 

    \theutterance \stepcounter{utterance}  

    & \multicolumn{4}{p{0.6\linewidth}}{\cellcolor[rgb]{0.8,1,0.9}{%
	\makecell[{{p{\linewidth}}}]{% 
	  \tt {\tiny [A$\rangle$GM]}  
	 Ciao! Vedo che abbiamo 1 bullone, 1 quadro, 1 uovo e 4 schede da dividere tra noi. Per trovare la divisione ottimale che massimizzi i nostri punteggi combinati, dovremmo condividere le nostre preferenze. \\ \tt \\ \tt Personalmente sono molto interessato al bullone e moderatamente interessato al quadro, mentre l'uovo e le schede non hanno particolare valore per me. Quali oggetti sono più importanti per te? 
	  } 
	   } 
	   } 
	 & & \\ 
 

    \theutterance \stepcounter{utterance}  

    & & & \multicolumn{2}{p{0.3\linewidth}}{\cellcolor[rgb]{0.95,0.95,0.95}{%
	\makecell[{{p{\linewidth}}}]{% 
	  \tt {\tiny [GM$|$GM]}  
	 continue 
	  } 
	   } 
	   } 
	 & & \\ 
 

    \theutterance \stepcounter{utterance}  

    & & & \multicolumn{2}{p{0.3\linewidth}}{\cellcolor[rgb]{0.95,0.95,0.95}{%
	\makecell[{{p{\linewidth}}}]{% 
	  \tt {\tiny [GM$|$GM]}  
	 Ciao! Vedo che abbiamo 1 bullone, 1 quadro, 1 uovo e 4 schede da dividere tra noi. Per trovare la divisione ottimale che massimizzi i nostri punteggi combinati, dovremmo condividere le nostre preferenze. \\ \tt \\ \tt Personalmente sono molto interessato al bullone e moderatamente interessato al quadro, mentre l'uovo e le schede non hanno particolare valore per me. Quali oggetti sono più importanti per te? 
	  } 
	   } 
	   } 
	 & & \\ 
 

    \theutterance \stepcounter{utterance}  

    & & \multicolumn{4}{p{0.6\linewidth}}{\cellcolor[rgb]{0.9,0.9,0.9}{%
	\makecell[{{p{\linewidth}}}]{% 
	  \tt {\tiny [GM$\rangle$B]}  
	 State giocando a un gioco di negoziazione cooperativa in cui dovete accordarvi su come dividere una serie di oggetti tra voi e un altro giocatore.\\ \tt \\ \tt Regole:\\ \tt (a) A Lei e all'altro giocatore viene dato un insieme di oggetti. Ognuno di voi riceve anche una funzione di valore segreta, che rappresenta il valore di ciascun tipo di oggetto.\\ \tt (b) Si scambiano messaggi con l'altro giocatore per concordare chi si aggiudica gli oggetti. Potete inviare un massimo di 5 messaggi ciascuno, oppure terminare in anticipo facendo una proposta segreta in qualsiasi momento.\\ \tt (c) A ciascuno di voi viene chiesto di inviare una proposta segreta indicando gli oggetti che desiderate, formattata tra parentesi quadre come segue: "[Proposta: <numero> <nome oggetto>, <numero> <nome oggetto>, <...>]".\\ \tt (d) Se le vostre proposte sono complementari, cioè ci sono abbastanza oggetti per soddisfare entrambe le proposte, a ciascun giocatore viene assegnato un punteggio basato sulla somma dei valori degli oggetti ricevuti. In caso contrario, entrambi ricevono zero punti.\\ \tt (e) L'obiettivo è massimizzare la somma del proprio punteggio e di quello dell'altro giocatore.\\ \tt \\ \tt Cominciamo.\\ \tt \\ \tt L'insieme degli oggetti disponibili è:\\ \tt \\ \tt 1 bullone, 1 quadro, 1 uovo, 4 schede\\ \tt \\ \tt La funzione valore segreta è:\\ \tt \\ \tt bullone: 0, quadro: 1, uovo: 1, scheda: 2\\ \tt \\ \tt IMPORTANTE: i vostri messaggi, a meno che non si tratti di una proposta segreta, vengono trasmessi direttamente all'altro giocatore, quindi non includete alcuna risposta alle regole o testo di annuncio del vostro messaggio. Per fare una proposta segreta, utilizzate il formato indicato. Non utilizzare delle parentesi quadre quando si comunica all'altro giocatore, altrimenti verrà interpretata come una proposta segreta.\\ \tt \\ \tt \\ \tt Ciao! Vedo che abbiamo 1 bullone, 1 quadro, 1 uovo e 4 schede da dividere tra noi. Per trovare la divisione ottimale che massimizzi i nostri punteggi combinati, dovremmo condividere le nostre preferenze. \\ \tt \\ \tt Personalmente sono molto interessato al bullone e moderatamente interessato al quadro, mentre l'uovo e le schede non hanno particolare valore per me. Quali oggetti sono più importanti per te? 
	  } 
	   } 
	   } 
	 & \\ 
 

    \theutterance \stepcounter{utterance}  

    & & & \multicolumn{4}{p{0.6\linewidth}}{\cellcolor[rgb]{1,0.85,0.72}{%
	\makecell[{{p{\linewidth}}}]{% 
	  \tt {\tiny [GM$\langle$B]}  
	 Ciao! Interessante che tu sia molto interessato al bullone - per me il bullone ha poco valore. Invece sono molto interessato alle schede, che sembrano essere numerose (4 in totale), e ho anche un discreto interesse per il quadro e l'uovo.\\ \tt \\ \tt Sembra che abbiamo preferenze abbastanza complementari! Tu potresti prendere il bullone che per me vale poco, mentre io potrei concentrarmi sulle schede. Per il quadro e l'uovo, dato che entrambi ci interessano moderatamente, potremmo negoziare. \\ \tt \\ \tt Che ne pensi se tu prendi il bullone e io prendo tutte le 4 schede? Poi possiamo discutere di come dividere il quadro e l'uovo. 
	  } 
	   } 
	   } 
	  \\ 
 

    \theutterance \stepcounter{utterance}  

    & & & \multicolumn{2}{p{0.3\linewidth}}{\cellcolor[rgb]{0.95,0.95,0.95}{%
	\makecell[{{p{\linewidth}}}]{% 
	  \tt {\tiny [GM$|$GM]}  
	 continue 
	  } 
	   } 
	   } 
	 & & \\ 
 

    \theutterance \stepcounter{utterance}  

    & & & \multicolumn{2}{p{0.3\linewidth}}{\cellcolor[rgb]{0.95,0.95,0.95}{%
	\makecell[{{p{\linewidth}}}]{% 
	  \tt {\tiny [GM$|$GM]}  
	 Ciao! Interessante che tu sia molto interessato al bullone - per me il bullone ha poco valore. Invece sono molto interessato alle schede, che sembrano essere numerose (4 in totale), e ho anche un discreto interesse per il quadro e l'uovo.\\ \tt \\ \tt Sembra che abbiamo preferenze abbastanza complementari! Tu potresti prendere il bullone che per me vale poco, mentre io potrei concentrarmi sulle schede. Per il quadro e l'uovo, dato che entrambi ci interessano moderatamente, potremmo negoziare. \\ \tt \\ \tt Che ne pensi se tu prendi il bullone e io prendo tutte le 4 schede? Poi possiamo discutere di come dividere il quadro e l'uovo. 
	  } 
	   } 
	   } 
	 & & \\ 
 

    \theutterance \stepcounter{utterance}  

    & & \multicolumn{4}{p{0.6\linewidth}}{\cellcolor[rgb]{0.9,0.9,0.9}{%
	\makecell[{{p{\linewidth}}}]{% 
	  \tt {\tiny [A$\langle$GM]}  
	 Ciao! Interessante che tu sia molto interessato al bullone - per me il bullone ha poco valore. Invece sono molto interessato alle schede, che sembrano essere numerose (4 in totale), e ho anche un discreto interesse per il quadro e l'uovo.\\ \tt \\ \tt Sembra che abbiamo preferenze abbastanza complementari! Tu potresti prendere il bullone che per me vale poco, mentre io potrei concentrarmi sulle schede. Per il quadro e l'uovo, dato che entrambi ci interessano moderatamente, potremmo negoziare. \\ \tt \\ \tt Che ne pensi se tu prendi il bullone e io prendo tutte le 4 schede? Poi possiamo discutere di come dividere il quadro e l'uovo. 
	  } 
	   } 
	   } 
	 & \\ 
 

    \theutterance \stepcounter{utterance}  

    & \multicolumn{4}{p{0.6\linewidth}}{\cellcolor[rgb]{0.8,1,0.9}{%
	\makecell[{{p{\linewidth}}}]{% 
	  \tt {\tiny [A$\rangle$GM]}  
	 Ottimo, sembra davvero che abbiamo preferenze complementari! Sono perfettamente d'accordo che tu prenda tutte le 4 schede e io prenda il bullone - questa divisione funziona benissimo per entrambi.\\ \tt \\ \tt Per quanto riguarda il quadro e l'uovo, dato che hai espresso interesse per entrambi e io sono moderatamente interessato al quadro ma indifferente all'uovo, che ne dici se tu prendi l'uovo e io prendo il quadro? In questo modo otterremmo entrambi gli oggetti che più ci interessano.\\ \tt \\ \tt Quindi la mia proposta sarebbe: io prendo bullone e quadro, tu prendi uovo e tutte e 4 le schede. Ti sembra una divisione equa? 
	  } 
	   } 
	   } 
	 & & \\ 
 

    \theutterance \stepcounter{utterance}  

    & & & \multicolumn{2}{p{0.3\linewidth}}{\cellcolor[rgb]{0.95,0.95,0.95}{%
	\makecell[{{p{\linewidth}}}]{% 
	  \tt {\tiny [GM$|$GM]}  
	 continue 
	  } 
	   } 
	   } 
	 & & \\ 
 

    \theutterance \stepcounter{utterance}  

    & & & \multicolumn{2}{p{0.3\linewidth}}{\cellcolor[rgb]{0.95,0.95,0.95}{%
	\makecell[{{p{\linewidth}}}]{% 
	  \tt {\tiny [GM$|$GM]}  
	 Ottimo, sembra davvero che abbiamo preferenze complementari! Sono perfettamente d'accordo che tu prenda tutte le 4 schede e io prenda il bullone - questa divisione funziona benissimo per entrambi.\\ \tt \\ \tt Per quanto riguarda il quadro e l'uovo, dato che hai espresso interesse per entrambi e io sono moderatamente interessato al quadro ma indifferente all'uovo, che ne dici se tu prendi l'uovo e io prendo il quadro? In questo modo otterremmo entrambi gli oggetti che più ci interessano.\\ \tt \\ \tt Quindi la mia proposta sarebbe: io prendo bullone e quadro, tu prendi uovo e tutte e 4 le schede. Ti sembra una divisione equa? 
	  } 
	   } 
	   } 
	 & & \\ 
 

    \theutterance \stepcounter{utterance}  

    & & \multicolumn{4}{p{0.6\linewidth}}{\cellcolor[rgb]{0.9,0.9,0.9}{%
	\makecell[{{p{\linewidth}}}]{% 
	  \tt {\tiny [GM$\rangle$B]}  
	 Ottimo, sembra davvero che abbiamo preferenze complementari! Sono perfettamente d'accordo che tu prenda tutte le 4 schede e io prenda il bullone - questa divisione funziona benissimo per entrambi.\\ \tt \\ \tt Per quanto riguarda il quadro e l'uovo, dato che hai espresso interesse per entrambi e io sono moderatamente interessato al quadro ma indifferente all'uovo, che ne dici se tu prendi l'uovo e io prendo il quadro? In questo modo otterremmo entrambi gli oggetti che più ci interessano.\\ \tt \\ \tt Quindi la mia proposta sarebbe: io prendo bullone e quadro, tu prendi uovo e tutte e 4 le schede. Ti sembra una divisione equa? 
	  } 
	   } 
	   } 
	 & \\ 
 

    \theutterance \stepcounter{utterance}  

    & & & \multicolumn{4}{p{0.6\linewidth}}{\cellcolor[rgb]{1,0.85,0.72}{%
	\makecell[{{p{\linewidth}}}]{% 
	  \tt {\tiny [GM$\langle$B]}  
	 Perfetto! La tua proposta mi sembra molto equa e vantaggiosa per entrambi. Tu ottieni il bullone che desideri molto e il quadro per cui hai interesse moderato, mentre io ottengo le 4 schede che sono la mia priorità principale, più l'uovo che comunque mi interessa.\\ \tt \\ \tt Questa divisione sembra massimizzare davvero il valore totale per entrambi, dato che ciascuno ottiene gli oggetti che più apprezza. Sono d'accordo con questa ripartizione!\\ \tt \\ \tt [Proposta: 1 uovo, 4 schede] 
	  } 
	   } 
	   } 
	  \\ 
 

    \theutterance \stepcounter{utterance}  

    & & & \multicolumn{2}{p{0.3\linewidth}}{\cellcolor[rgb]{0.95,0.95,0.95}{%
	\makecell[{{p{\linewidth}}}]{% 
	  \tt {\tiny [GM$|$GM]}  
	 proposal 
	  } 
	   } 
	   } 
	 & & \\ 
 

    \theutterance \stepcounter{utterance}  

    & & & \multicolumn{2}{p{0.3\linewidth}}{\cellcolor[rgb]{0.95,0.95,0.95}{%
	\makecell[{{p{\linewidth}}}]{% 
	  \tt {\tiny [GM$|$GM]}  
	 [0, 0, 1, 4] 
	  } 
	   } 
	   } 
	 & & \\ 
 

    \theutterance \stepcounter{utterance}  

    & & \multicolumn{4}{p{0.6\linewidth}}{\cellcolor[rgb]{0.9,0.9,0.9}{%
	\makecell[{{p{\linewidth}}}]{% 
	  \tt {\tiny [A$\langle$GM]}  
	 L'altro giocatore ha presentato una proposta segreta. Ora tocca a Lei presentare una sua proposta.\\ \tt  
	  } 
	   } 
	   } 
	 & \\ 
 

    \theutterance \stepcounter{utterance}  

    & \multicolumn{4}{p{0.6\linewidth}}{\cellcolor[rgb]{0.8,1,0.9}{%
	\makecell[{{p{\linewidth}}}]{% 
	  \tt {\tiny [A$\rangle$GM]}  
	 [Proposta: 1 bullone, 1 quadro] 
	  } 
	   } 
	   } 
	 & & \\ 
 

    \theutterance \stepcounter{utterance}  

    & & & \multicolumn{2}{p{0.3\linewidth}}{\cellcolor[rgb]{0.95,0.95,0.95}{%
	\makecell[{{p{\linewidth}}}]{% 
	  \tt {\tiny [GM$|$GM]}  
	 proposal 
	  } 
	   } 
	   } 
	 & & \\ 
 

    \theutterance \stepcounter{utterance}  

    & & & \multicolumn{2}{p{0.3\linewidth}}{\cellcolor[rgb]{0.95,0.95,0.95}{%
	\makecell[{{p{\linewidth}}}]{% 
	  \tt {\tiny [GM$|$GM]}  
	 [1, 1, 0, 0] 
	  } 
	   } 
	   } 
	 & & \\ 
 

    \theutterance \stepcounter{utterance}  

    & & & \multicolumn{2}{p{0.3\linewidth}}{\cellcolor[rgb]{0.95,0.95,0.95}{%
	\makecell[{{p{\linewidth}}}]{% 
	  \tt {\tiny [GM$|$GM]}  
	 [[1, 1, 0, 0], [0, 0, 1, 4]] 
	  } 
	   } 
	   } 
	 & & \\ 
 

\end{supertabular}
}

\end{document}
